\documentclass[../../main.tex]{subfiles}
\graphicspath{{\subfix{../../image/}}}

\begin{document}

\section{样本空间与事件}\label{section:ch01-样本空间与事件}

\subsection*{一、样本空间}

从本节开始，我们将逐步引进概率论的基本概念。样本空间与事件是最基本的两个概念。

对随机现象的研究必然要联系到对客观的事物进行“调查”、“观察”或“实验”，以后我们统称之为（随机）\textbf{试验}（trial），并假定这种“试验”可以在相同条件下重复进行。

我们感兴趣的是试验的结果。例如掷一次硬币，我们关心的是出现正面或出现反面，这是两个可能出现的结果。假如我们考察的是掷两次硬币的试验，则可能出现的结果有（正，正），（正，反），（反，正），（反，反）四种；如果掷三次硬币，则结果还要复杂，但还是可以把它们描述出来。总之，为了研究随机试验，首先需要知道这个试验可能出现的结果。这些结果称为\textbf{样本点}（sample point），一般用 $\omega$ 表示。样本点全体构成\textbf{样本空间}（sample space），用 $\Omega$ 表示。在具体问题中，给定样本空间是描述随机现象的第一步。

下面举一些例子。

\begin{example}[英文字母]\label{example:ch01-英文字母}
在研究英文字母使用情况时，把样本空间选为 $\Omega=\{\text{空格},A,B,\cdots,Z\}$ 是适宜的，这个样本空间只有有限个样本点，是比较简单的样本空间。
\end{example}

\begin{example}[粒子数]\label{example:ch01-粒子数}
观察一小时中落在地球上某一区域的粒子数，可能的结果一定是非负整数，而且很难指定一个数作为它的上限，这样，可以把样本空间取为 $\Omega=\{0,1,2,\cdots\}$。这个样本空间含有无穷多个样本点，但这些样本点可以依照某种次序排列出来，以后我们将称它的点数为可列个。
\end{example}

\begin{example}[气温]\label{example:ch01-气温}
讨论某地区的气温时，我们自然把样本空间取为 $\Omega=(-\infty,\infty)$，或 $\Omega=[a,b]$。这个样本空间包含有无穷多个样本点，它们充满一个区间，不是一个可列集。
\end{example}

\begin{example}[震源]\label{example:ch01-震源}
考察地震震源时，可以把样本点取为 $(x,y,z)$，其中 $x$ 表示震源的经度，$y$ 表示纬度，$z$ 表示深度。这时，样本空间是三维空间中某一区域。
\end{example}

\begin{example}[道·琼斯指数]\label{example:ch01-道琼斯指数}
金融分析师把道·琼斯指数作为研究对象，每日的指数涨跌用一条曲线 $x(t),0\leqslant t\leqslant T$ 表示，作为一个样本点，这时样本空间是函数空间，这类样本空间是随机过程（stochastic process）理论的研究对象。
\end{example}

从以上例子可以看出，随着问题不同，样本空间可以相当简单，也可以十分复杂。

在今后讨论中，经常把样本空间认为是预先给定的。当然对于一个实际问题或一个随机现象，如何用一个恰当的样本空间来描述它也很值得研究。但是在概率论的研究中，一般都假定样本空间是给定的。这是必要的抽象，这种抽象使我们能更好地把握住随机现象的本质，而且得到的结果能广泛地应用。事实上，一个样本空间可以概括各种实际内容很不相同的问题：例如只包含两个样本点的样本空间既能作为掷硬币出现正、反面的模型，也能用于产品检验中出现“合格品”及“废品”，又能用于气象中“下雨”与“不下雨”，以及公用事业排队现象中“有人排队”与“无人排队”等等。尽管问题的实际内容如此不同，但有时却能归结为相同的概率模型。我们后面常以摸球等作为例子也是由于这个原因，它能使问题的本质更为突出。

\subsection*{二、事件}

有了样本空间的概念，就可以定义事件。我们还是从考察一个例子开始。

\begin{example}[摸球问题]\label{example:ch01-摸球}
口袋中装有 $4$ 只白球和 $2$ 只黑球，我们考虑依次从中摸出两球所可能出现的事件。若对球进行编号，$4$ 只白球分别编为 $1,2,3,4$ 号，$2$ 只黑球编为 $5,6$ 号。如果用数对 $(i,j)$ 表示第一次摸得 $i$ 号球，第二次摸得 $j$ 号球，则可能出现的结果是
\begin{align*}
&(1,2),(1,3),(1,4),(1,5),(1,6)\\
&(2,1),(2,3),(2,4),(2,5),(2,6)\\
&(3,1),(3,2),(3,4),(3,5),(3,6)\\
&(4,1),(4,2),(4,3),(4,5),(4,6)\\
&(5,1),(5,2),(5,3),(5,4),(5,6)\\
&(6,1),(6,2),(6,3),(6,4),(6,5)
\end{align*}
把这 $30$ 个结果作为样本点，则构成了样本空间。在这个问题中，这些样本点是我们感兴趣的事件；但是我们也可以研究下面另外一些事件：

$A$：第一次摸出黑球；

$B$：第二次摸出黑球；

$C$：第一次及第二次都摸出黑球。

后面这些事件与前面那些事件的不同处在于这些事件是可以分解的，例如为了 $A$ 出现必须而且只需下列样本点之一出现：
\begin{align*}
&(5,1),(5,2),(5,3),(5,4),(5,6)\\
&(6,1),(6,2),(6,3),(6,4),(6,5)
\end{align*}
前面的 $30$ 个事件由单个样本点构成；后面这三个事件，每一个事件都是由若干个样本点构成的，总之，它们都是样本点的某个集合。
\end{example}

所谓给定一个点的集合 $S$，是指对于任何一个点 $\omega$，都可以确定它是不是属于 $S$。如果是，则记为 $\omega\in S$；如果不是，则记为 $\omega\overline{\in}S$。按照这种定义，单个点也是一个点集。习惯上还约定不包含任何点的集合也是一个点集，称为空集，记为 $\varnothing$。

今后，我们把事件定义为样本点的某个集合，称某事件发生当且仅当它所包含的某一个样本点出现。

因此，虽然试验的全部可能结果在试验前就很明确，但是只有到了试验之后，才能确定某一特定的事件是否发生。

我们把样本空间 $\Omega$ 也作为一个事件，因为在每次试验中必然出现 $\Omega$ 中的某个样本点，也即 $\Omega$ 必然发生，所以常称 $\Omega$ 为\textbf{必然事件}。类似地，我们把空集 $\varnothing$ 也作为一个事件，它在每次试验中都不会发生，称为\textbf{不可能事件}。

必然事件 $\Omega$ 在试验中必然发生。相反地，不可能事件 $\varnothing$ 在任何试验中不可能发生，必然事件与不可能事件可以说不是随机事件，但为了今后研究的方便，我们还是把必然事件与不可能事件作为随机事件的两个极端情形来统一处理。

\subsection*{三、事件的运算}

在一个样本空间中显然可以定义不止一个事件。概率论的重要研究课题之一是希望从简单事件的概率推算出复杂事件的概率。在实际生活中，往往要求我们同时考察几个在同样条件下的事件以及它们之间的联系。详细地分析事件之间的关系，不仅帮助我们更深刻地认识事件的本质，而且可以大大简化一些复杂事件的概率计算。

下面就讨论事件间的关系及事件的运算，先讨论两个事件 $A$ 与 $B$ 之间的关系。

若 $A$ 中的每一个样本点都包含在 $B$ 中，则记为 $A\subset B$ 或 $B\supset A$，并称 $A$ 被包含于 $B$，亦称事件 $B$ 包含了事件 $A$，这时事件 $A$ 发生必然导致事件 $B$ 发生。例如若以 $A$ 记“来到呼叫不超过 $5$ 个”，以 $B$ 记“来到呼叫不超过 $6$ 个”，则 $A\subset B$。显然对任何事件 $A$，必有 $\varnothing\subset A\subset\Omega$。

如果 $A\supset B$ 与 $B\supset A$ 同时成立，则称 $A$ 与 $B$ \textbf{等价}或称 $A$ \textbf{等于} $B$，记为 $A=B$，等价的两个事件同时发生，因此可看作是一样的。

对于事件 $A$，由所有不包含在 $A$ 中的样本点所组成的事件称为 $A$ 的\textbf{逆事件}，或称 $A$ 的\textbf{对立事件}，记为 $\overline{A}$，$\overline{A}$ 表示 $A$ 不发生。例如若以 $A$ 表示“来到呼叫不超过 $5$ 个”，则 $\overline{A}$ 表示“来到呼叫超过 $5$ 个”。显然，若 $\overline{A}$ 是 $A$ 的对立事件，则 $A$ 也是 $\overline{A}$ 的对立事件，即 $\overline{\overline{A}}=A$。必然事件与不可能事件互为对立事件。

其次，对于事件 $A$ 及事件 $B$，定义两个新事件：

用 $A\cap B$ 或 $AB$ 表示所有同时属于 $A$ 及 $B$ 的样本点的集合，称它为 $A$ 与 $B$ 的\textbf{交}，事件 $AB$ 表示事件 $A$ 与事件 $B$ 同时发生。

用 $A\cup B$ 表示至少属于 $A$ 或 $B$ 中的一个的所有样本点的集合，称它为 $A$ 与 $B$ 的\textbf{并}，事件 $A\cup B$ 表示事件 $A$ 或事件 $B$ 或它们二者发生，也即表示事件 $A$ 与事件 $B$ 至少发生一个。

若 $AB=\varnothing$，则表示 $A$ 与 $B$ 不可能同时发生，称 $A$ 与 $B$ \textbf{互不相容}。样本点是互不相容的。

本书特别约定：对于互不相容事件 $A$ 与 $B$，我们称它们的并为\textbf{和}，并记作 $A+B$。

用 $A-B$ 表示包含在 $A$ 中而不包含在 $B$ 中的样本点全体，称之为 $A$ 与 $B$ 的\textbf{差}，事件 $A-B$ 表示事件 $A$ 发生而事件 $B$ 不发生。

关于事件运算的顺序作如下约定：先进行逆的运算，再进行交的运算，最后才进行并或差的运算。这与数式运算中，先函数，再乘除，后加减的约定相似。

用上面的记号可以把对立事件之间的关系表述如下：$A\cup\overline{A}=\Omega$，$A\cap\overline{A}=\varnothing$，这也可以作为对立事件的定义。显然 $\overline{A}=\Omega-A$。

事件运算成立如下\textbf{De Morgan 定理}，亦称对偶原理：
\begin{align}
\overline{A\cup B}=\overline{A}\cap\overline{B},\qquad \overline{A\cap B}=\overline{A}\cup\overline{B}\label{align:ch01-de-morgan}
\end{align}

于是，定义了逆运算后，交与并可以互相表示，差也可以用它们表示：
\begin{align*}
A\cup B=\overline{\overline{A}\cap\overline{B}},\qquad A\cap B=\overline{\overline{A}\cup\overline{B}}
\end{align*}
\begin{align*}
A-B=A\cap\overline{B}=\overline{\overline{A}\cup B}
\end{align*}

因此本质上只需要两种运算：“交与逆”或“并与逆”，不过定义 $4$ 种运算自有其方便之处。

有时用平面上某正方形中的图形来表示事件间的关系或运算较为直观，这种表示法称为\textbf{Venn 图}。

事件 $\Omega,A,\overline{A},A\cup B,AB,A-B$ 在图~\ref{figure:ch01-事件运算} 中分别以阴影表出。不难理解，$B\supset A$ 相应于 $A$ 的图形完全包含在 $B$ 的图形中；$A$ 和 $B$ 互不相容，则相应于 $A$ 和 $B$ 的图形不相交。

\begin{figure}[H]
\centering
\includegraphics[width=0.85\textwidth]{ch01-venn.png}
\caption{事件运算}\label{figure:ch01-事件运算}
\end{figure}

我们用例~\ref{example:ch01-摸球} 来说明这些关系。前面 $30$ 种可能结果就是样本点全体，它们构成必然事件 $\Omega$。

$A$ 记第一次摸得黑球，则它由第 $5$ 行及第 $6$ 行的 $10$ 个样本点构成；这时 $\overline{A}$ 表示第一次摸得白球，它由第一行至第四行的 $20$ 个样本点构成。显然 $A$ 与 $\overline{A}$ 互不相容，而且 $A+\overline{A}=\Omega$。

$B$ 记第二次摸得黑球，它由下列 $10$ 个样本点构成：
\begin{align*}
&(1,5),(1,6),(2,5),(2,6),(3,5)\\
&(3,6),(4,5),(4,6),(5,6),(6,5)
\end{align*}

事件 $A\cup B$ 表示第一次或第二次中至少有一次摸得黑球，它包含下列 $18$ 个样本点：
\begin{align*}
&(1,5),(1,6),(2,5),(2,6),(3,5),(3,6)\\
&(4,5),(4,6),(5,1),(5,2),(5,3),(5,4)\\
&(5,6),(6,1),(6,2),(6,3),(6,4),(6,5)
\end{align*}

事件 $AB$ 表示两次都摸得黑球，它由下列两个样本点构成：$(5,6),(6,5)$，这是 $A$ 与 $B$ 共同包含的样本点。因此 $C=AB$。

事件 $A-B$ 表示第一次摸得黑球而第二次摸出白球，包含了 $(5,1),(5,2),(5,3),(5,4),(6,1),(6,2),(6,3),(6,4)$ 等 $8$ 个样本点。

不难把上面定义推广到多个事件的场合。

例如，对于 $n$ 个事件 $A_1,A_2,\cdots,A_n$，用 $A_1\cup A_2\cup\cdots\cup A_n$ 或 $\bigcup_{i=1}^{n}A_i$ 表示 $A_1,A_2,\cdots,A_n$ 中至少发生一个，称为 $A_1,A_2,\cdots,A_n$ 的\textbf{并}，特别当 $A_1,A_2,\cdots,A_n$ 两两互不相容时，并特称为\textbf{和}，记作 $A_1+A_2+\cdots+A_n$ 或 $\sum_{i=1}^{n}A_i$。相应地，用 $A_1A_2\cdots A_n$ 或 $\bigcap_{i=1}^{n}A_i$ 表示 $A_1,A_2,\cdots,A_n$ 同时发生等等。

这时，对偶原理仍然成立：
\begin{align*}
\overline{\bigcup_{i=1}^{n}A_i}=\bigcap_{i=1}^{n}\overline{A_i},\qquad \overline{\bigcap_{i=1}^{n}A_i}=\bigcup_{i=1}^{n}\overline{A_i}
\end{align*}
并且有很明显的概率意义：至少发生一个的对立面是一个也不发生；全部发生的对立面是至少有一个不发生。

对于可列个事件的场合，我们定义
\begin{align*}
\bigcup_{i=1}^{\infty}A_i=\lim_{n\to\infty}\bigcup_{i=1}^{n}A_i,\qquad \bigcap_{i=1}^{\infty}A_i=\lim_{n\to\infty}\bigcap_{i=1}^{n}A_i
\end{align*}
这时依然有对偶原理。

事件的运算成立下列关系式，它们的证明留给读者。

（1）\textbf{交换律}：$A\cup B=B\cup A$，$AB=BA$；

（2）\textbf{结合律}：$(A\cup B)\cup C=A\cup(B\cup C)$，$(AB)C=A(BC)$；

（3）\textbf{分配律}：$(A\cup B)\cap C=AC\cup BC$，$(A\cap B)\cup C=(A\cup C)\cap(B\cup C)$。

熟悉集合论的读者或许早就发现，事件间的关系及运算与集合论中或 Boole（$1815$—$1864$）代数中集合的关系及运算是完全相似的，而且这个相似性在建立概率论的严格数学基础时非常重要。不过，我们应该强调另一面，就是要学会用概率论的语言来解释这些关系及运算，并且会用这些运算关系来表示各种事件。

\begin{example}[三事件关系]\label{example:ch01-三事件关系}
若 $A,B,C$ 是三个事件，则

（1）所有这三个事件都发生可以表示为：$ABC$；

（2）这三个事件恰好发生一个可以表示为：$A\overline{B}\overline{C}+\overline{A}B\overline{C}+\overline{A}\overline{B}C$；

（3）这三个事件恰好发生两个可以表示为：$AB\overline{C}+A\overline{B}C+\overline{A}BC$；

（4）这三个事件中至少发生一个可以表示为：$A\cup B\cup C$ 或 $A\overline{B}\overline{C}+\overline{A}B\overline{C}+\overline{A}\overline{B}C+AB\overline{C}+A\overline{B}C+\overline{A}BC+ABC$；

还有一种看似复杂的表示法：$\overline{\overline{A}\overline{B}\overline{C}}$，正是对偶公式，今后很有用。

（5）$A$ 发生而 $B$ 与 $C$ 都不发生可以表示为：$A\overline{B}\overline{C}$ 或 $A-B-C$ 或 $A-(B\cup C)$；

（6）$A$ 与 $B$ 都发生而 $C$ 不发生可以表示为：$AB\overline{C}$ 或 $AB-C$ 或 $AB-ABC$。
\end{example}

\subsection*{四、有限样本空间}

我们先考虑只有有限个样本点的样本空间，这种样本空间称为\textbf{有限样本空间}。这是最简单的样本空间，研究它有助于深入研究更为复杂的样本空间。

若 $\Omega$ 是有限样本空间，其样本点为 $\omega_1,\omega_2,\cdots,\omega_n$，在这种场合可以把 $\Omega$ 的任何子集都当作事件。这时，样本点作为单点集，当然是事件。在这种样本空间中引进概率，只要对每个样本点 $\omega_i$，给定一个数与它对应，此数称为 $\omega_i$ 的概率，并记之为 $P(\omega_i)$，它是非负的，而且满足
\begin{align}
P(\omega_1)+P(\omega_2)+\cdots+P(\omega_n)=1\label{align:ch01-有限样本空间概率和}
\end{align}
这样，我们对样本点定义了概率，用它来度量每个样本点出现的可能性的大小。由此出发，我们不难定义更为一般的事件的概率。

\begin{definition}[事件的概率]\label{definition:ch01-事件概率}
任何事件 $A$ 的概率 $P(A)$ 是 $A$ 中各样本点的概率之和。
\end{definition}

按照这个定义，显然有 $P(\Omega)=1$，$0\leqslant P(A)\leqslant 1$。

如在例~\ref{example:ch01-摸球} 中，若定义每个样本点出现的概率均为 $\frac{1}{30}$（这相当于假定各个球外形完全一样，并且摸球是随机的，各个球被摸到的机会均等），则得 $P(A)=\frac{10}{30}$，$P(B)=\frac{10}{30}$，$P(C)=\frac{2}{30}$，$P(A\cup B)=\frac{18}{30}$，$P(A-B)=\frac{8}{30}$ 等等。

我们将在下一节研究一种特殊的有限样本空间。

把上面做法推广到有可列个样本点的样本空间是不难的，这种空间称为\textbf{离散样本空间}。但是当把上面做法推广到不可列个样本点的场合，则会遇到实质性的困难，对于这种一般场合的讨论，以后将逐渐展开。

\end{document}